\documentclass{article}

\usepackage{tabularx}
\usepackage{booktabs}

\title{Problem Statement and Goals\\\progname}

\author{\authname}

\date{}

\input{../Comments.text}
\input{../Common.text}

\begin{document}

\maketitle

\begin{table}[hp]
\caption{Revision History} \label{TblRevisionHistory}
\begin{tabularx}{\textwidth}{llX}
\toprule
\textbf{Date} & \textbf{Developer(s)} & \textbf{Change}\\
\midrule
September 24, 2026 & Andreas Sideris & Problem Statement added\\
September 25, 2026 & Andreas Sideris & Inputs and Outputs completed\\
... & ... & ...\\
\bottomrule
\end{tabularx}
\end{table}

\section{Problem Statement} 

\subsection{Problem}

Video games of the roguelite genre have skyrocketed in popularity in recent 
years, with titles such as \textit{Hades}, \textit{Slay the Spire}, and 
\textit{Balatro} releasing to widespread acclaim. This has resulted in the 
development of many new projects aiming to follow in their footsteps, 
often oversaturating the genre with redundant, unengaging gameplay in certain 
cases. 
Due to the 
unique nature of roguelite video games (users repeat `runs' of the game after 
their 
progress resets from each failed game session, and attempt to upgrade their 
ability 
to succeed between each 
run), poorly designed titles can significantly negatively affect their user 
experience. This can manifest due to a failure to conduct effective, 
quantifiable 
usability 
testing, to ensure that a game's roguelite features are engaging 
and satisfying to play. To address this issue, this project will involve the 
development of a small, unique roguelite video game that provides explicit 
variation 
capabilities for its user interface and gameplay, to facilitate more impactful 
usability 
testing.

\subsection{Inputs and Outputs}

\textbf{Inputs:}

\begin{enumerate}
	\item Users' gameplay control inputs
	\item Usability metric data collected from playtests involving feature 
	variations
	\item Individual feedback collected from playtesters
\end{enumerate}
	
\noindent \textbf{Outputs:}

\begin{enumerate}
	\item Using input 1: Gameplay display outputs (user interface, characters 
	and objects, action feedback)
	\item Using inputs 2 and 3: Refined gameplay and user interface features
\end{enumerate}

\subsection{Stakeholders}

\subsection{Environment}

\wss{Hardware and Software Environment for the users.  The developer environment
is summarized as part of the developer plan.}

\section{Goals}

\section{Stretch Goals}

\wss{Goals you hope to get to, but not necessary.  They are also helpful if the
scope of the original project needs to be expanded.}

\section{Custom Documents (CDs)}

\wss{For CAS 741: State whether the project is a research project. This
designation, with the approval (or request) of the instructor, can be modified
over the course of the term.}

\wss{For SE Capstone: List your custom documents (CDs).  Potential CDs include
usability testing, code walkthroughs, user documentation, formal proof,
GenderMag personas, Design Thinking, etc.  (The full list is on the course
outline and in Lecture 02.) Normally the number of CDs will be two (one per
term).  Approval of the CDs will be part of the discussion with the instructor
for approving the project. The CDs, with the approval (or request) of the
instructor, can be modified over the course of the term.}

\newpage{}

\section{Appendix --- Generative AI Use Disclosure}

\wss{ChatGPT (version 5) was used to brainstorm ideas for the template for this 
disclose section.}

\wss{This section is about the use of AI to write this document, not your
	planned use of AI for your project.  You discuss that topic in one of the
	sections in the body of this report.}

\subsection{AI Tools Used}

\wss{Give the name of the AI tool(s) used and the version.}

\subsection{How and Where Used}

\wss{Explain what the tool(s) were used for.  Examples include: brainstorming,
	explaining a technical concept, reviewing the document, generating or 
	modifying
	text, generating or modifying diagrams, identifying errors or omissions or
	inconsistencies.}

\wss{For each of the uses, identify which parts of the document were affected.}

\wss{You do not need to report every interaction, but you should disclose the 
	uses that materially contributed to your work}

\subsection{Verification of AI-Generated Content}

\wss{\begin{itemize}
		\item Describe how the team checked AI-generated or AI-assisted material
		before including it in the document.
		\item In particular, verify factual claims, technical assertions,
		requirements, calculations, references, and other information for which 
		an
		AI system may produce incorrect or fabricated results.
		\item \textbf{The team is responsible for the accuracy, completeness,
			consistency, and appropriateness of the submitted document, 
			regardless of
			whether material was generated by a human or an AI system.}
\end{itemize}}

\newpage{}

\section*{Appendix --- Reflection}

\wss{Not required for CAS 741}

\input{../Reflection.text}

\begin{enumerate}
    \item What went well while writing this deliverable? 
    \item What pain points did you experience during this deliverable, and how
    did you resolve them?
    \item How did you and your team adjust the scope of your goals to ensure
    they are suitable for a Capstone project (not overly ambitious but also of
    appropriate complexity for a senior design project)?
\end{enumerate}  

\end{document}